\chapter{Durchführung} \label{ch:Durchfuehrung}

\section{Versuchsaufbau} \label{sec:aufbau}
\img{Versuchsaufbau}{Versuchsaufbau mit Beschriftung. \cite{skriptPAP22}}

\section{Messverfahren} \label{sec:Messverfahren}
Der Versuch wurde mit einem PC-gesteuerten Funktionsgenerator als Signalquelle und einem Analog-Oszilloskop zur Visualisierung und Messung der Spannungsverläufe durchgeführt. Als passive Bauelemente dienten verschiedene Widerstände, Kondensatoren und eine Spule, die auf einem Steckbrett zu den jeweiligen Schaltungen montiert wurden. Zusätzlich kamen ein Impedanzwandler, ein Niederfrequenz-Verstärker sowie eine Ladungsdrahtantenne zum Einsatz.

Zur Bestimmung der Zeitkonstante von RC-Gliedern wurde eine Rechteckspannung an die Schaltung angelegt. Die Halbwertszeit $T_{1/2}$ wurde direkt am Oszilloskop abgelesen, indem der Zeitpunkt bestimmt wurde, an dem die Spannung am Kondensator (bzw. Widerstand) die Hälfte des Endwertes erreicht hatte. Aus diesem Wert wurde die experimentelle Zeitkonstante gemäß \ceq{rc_halbwertszeit} berechnet und mit dem theoretischen Wert $\tau = RC$ verglichen.

Die Untersuchung des RC-Glieds als Integrator und Differentiator erfolgte qualitativ durch die Beobachtung der Ausgangsspannung bei verschiedenen Eingangssignalformen und Variation der Zeitkonstante durch Änderung des Widerstandswertes.

Für die Analyse des Frequenz- und Phasengangs von Hoch- und Tiefpassfiltern wurde ein Circuit Analyzer verwendet. Die Grenzfrequenz wurde sowohl grafisch aus dem Amplitudengang (Abfall auf $1/\sqrt{2}$) als auch über den Phasengang (Phasenverschiebung von $45^\circ$) bestimmt.

Der Serienschwingkreis wurde durch Variation des Vorwiderstands $R$ untersucht. Die Resonanzfrequenz wurde durch Abtasten des Frequenzbereichs und Identifikation des Maximums der Stromamplitude (bzw. Spannung am Widerstand) ermittelt. Die Bandbreite $\Delta f$ wurde als Frequenzdifferenz der Punkte bestimmt, an denen die Amplitude auf das $1/\sqrt{2}$-fache des Maximalwertes abgefallen ist. Daraus wurden die Induktivität der Spule sowie der Verlustwiderstand berechnet. Zusätzlich wurde der Verlustwiderstand über die Spannungsteilerregel im Resonanzfall bestimmt, indem das Verhältnis von Ein- und Ausgangsspannung gemessen wurde.

Die Bestimmung der Dämpfungskonstante eines freien, gedämpften Schwingkreises erfolgte durch Anregung mit einem kurzen Impuls und anschließende Beobachtung des freien Abklingverhaltens. Mehrere aufeinanderfolgende Amplitudenmaxima wurden am Oszilloskop abgelesen, um das logarithmische Dekrement $\Lambda$ zu berechnen.

Zur Untersuchung der Resonanzüberhöhung wurden die Resonanzfrequenzen separat für den Abgriff über den Widerstand, den Kondensator und die Spule bestimmt. Dabei wurde die Frequenz variiert, bis das Maximum der jeweiligen Teilspannung erreicht war.

Der Parallelschwingkreis wurde als Bandsperre aufgebaut. Die Resonanzfrequenz wurde durch Suche nach dem Minimum der Ausgangsspannung (bzw. Maximum der Impedanz) ermittelt.

Abschließend wurde ein Signal, bestehend aus drei überlagerten Sinusfrequenzen, durch verschiedene Filter (Hochpass, Tiefpass, Bandpass mit unterschiedlichen Widerstandswerten) geleitet. Die Dämpfung der einzelnen Frequenzanteile wurde durch Vergleich der Amplituden vor und nach dem Filter bestimmt, wobei die Spannungen in Volt umgerechnet wurden.
